%\chapter{\MakeUppercase{Recomendaciones}}\thispagestyle{fancy}
\chapter{Recomendaciones}\thispagestyle{fancy}
%\addcontentsline{toc}{chapter}{Recomendaciones}
\begin{itemize}
	\item Es preciso desarrollar el análisis teórico del error que valide las experiencias numéricas presentadas en este trabajo.
	\item Es posible resolver la ecuación de evolución  en una región muy próxima a la frontera libre pues solo interesa el comportamiento de una curva de nivel y no es necesario resolver sobre todo el dominio rectangular, este enfoque se conoce en la literatura como \emph{narrow band level set method}.
	\item Es importante que la función conjunto de nivel  $\phi$ conserve la propiedad de distancia dirigida en una región próxima a la frontera libre $\Gamma$ mediante algún proceso de reinicialización que controle el valor de $\|\nabla \phi \|$.
	\item Es necesario considerar el proceso de necrosis, es decir si la concentración de nutrientes cae por debajo de un umbral dado entonces las células de tumor mueren, esto da lugar a problemas de Poisson con término fuente. 
\end{itemize}
